% =============================================================================
% symmetry_spmv_optimizations.tex
%
% Algorithmic speedup & implementation of the fixed-Sz + spatial-symmetry
% construction and the symmetrized matrix-vector product (QED, N=32 regime).
%
% Build: pdflatex symmetry_spmv_optimizations.tex
% =============================================================================
\documentclass[11pt]{article}

\usepackage[margin=1in]{geometry}
\usepackage{amsmath,amssymb}
\usepackage{booktabs}
\usepackage{array}
\usepackage{xcolor}
\usepackage{listings}
\usepackage[hidelinks]{hyperref}
\usepackage{microtype}

% ---- code listing style (C++ / CUDA) ---------------------------------------
\definecolor{codegray}{gray}{0.45}
\definecolor{codekw}{rgb}{0.10,0.10,0.60}
\definecolor{codestr}{rgb}{0.55,0.10,0.10}
\definecolor{codecmt}{rgb}{0.20,0.45,0.20}
\lstdefinestyle{cpp}{
  language=C++,
  basicstyle=\ttfamily\footnotesize,
  keywordstyle=\color{codekw}\bfseries,
  stringstyle=\color{codestr},
  commentstyle=\color{codecmt}\itshape,
  numberstyle=\tiny\color{codegray},
  numbers=left,
  numbersep=6pt,
  showstringspaces=false,
  breaklines=true,
  columns=fullflexible,
  frame=single,
  framesep=4pt,
  rulecolor=\color{codegray},
  morekeywords={uint64_t,int32_t,cuDoubleComplex,__device__,__global__,__constant__,
                __forceinline__,std,size_t,nodiscard,constexpr,noexcept,Complex}
}
\lstset{style=cpp}

\newcommand{\bigO}{\mathcal{O}}
\newcommand{\Norm}{\mathcal{N}}
\newcommand{\reps}{\texttt{reps}}
\newcommand{\code}[1]{\texttt{\small #1}}

\title{Algorithmic Speedup and Implementation of the\\
       Fixed-$S^z$ + Symmetry Construction and Symmetrized SpMV}
\author{QED exact-diagonalization library}
\date{June 2026}

\begin{document}
\maketitle

\begin{abstract}
We document the algorithm- and implementation-level optimizations that make a
32-site fixed-$S^z$ plus spatial-symmetry exact-diagonalization workload
(microcanonical TPQ) feasible on a single large GPU. Two hot paths are covered:
(i) the \emph{symmetry construction} (orbit-representative enumeration, sector
dimensioning, and lazy / CSR-free sector data), and (ii) the \emph{symmetrized
matrix--vector product} (SpMV). The central result is replacing an orbit-CSR
``expand $\rightarrow$ apply $\rightarrow$ project'' SpMV---whose
$\mathcal{O}(\text{full-}S^z\text{ dim})$ scatter/gather tables are streamed over
PCIe every iteration---with an \emph{on-the-fly representative} SpMV that
regenerates the group action arithmetically on the device, turning the
$|G|$-fold dimension reduction into a genuine $|G|$-fold matvec speedup while
dropping the host working set from ${\sim}24$\,GiB/sector to ${\sim}1.2$\,GiB.
\end{abstract}

\section{Notation and setting}

For $N$ sites at magnetization $n_{\uparrow}$, the fixed-$S^z$ subspace has
dimension
\[
  D \;=\; \binom{N}{n_{\uparrow}},
  \qquad
  D\big|_{N=32,\,n_{\uparrow}=16} = 601{,}080{,}390 \approx 6.0\times 10^{8}.
\]
A spatial symmetry group $G$ (here $|G|=8$ translations) acts on computational
basis states by site permutations $g$. It partitions the $D$ states into
\emph{orbits}; each orbit is labeled by its canonical \emph{representative}
\[
  r_\alpha \;=\; \min_{g\in G} g(s),\qquad s\in \text{orbit}(\alpha).
\]
A momentum/irrep sector $k$ keeps one symmetrized basis vector per surviving
orbit,
\begin{equation}
  |\alpha\rangle
   \;=\; \frac{1}{\Norm_\alpha}\sum_{g\in G}\chi_k(g)^{*}\, g\,|r_\alpha\rangle,
  \label{eq:symvec}
\end{equation}
with characters $\chi_k(g)$ and normalization $\Norm_\alpha$. The sector
dimension is $\approx D/|G| \approx 7.5\times 10^{7}$. The goal is to build and
apply $H$ in this reduced basis \emph{without} ever paying $\bigO(D)$ memory.

Throughout, ``orbit CSR'' denotes the compressed structure storing, per
representative, all $|G|$ computational images $\{g\,|r_\alpha\rangle\}$ and
their coefficients $\chi_k(g)^{*}$---an $\bigO(D)$ object (${\approx}14$\,GiB at
$N{=}32$) plus an $\bigO(D)$ reverse index (${\approx}10$\,GiB).

\section{Symmetry construction}

\subsection{Orbit-representative enumeration: hash-set dedup $\rightarrow$
            min-image identity}

The naive enumerator scans all $D$ states and de-duplicates representatives in a
hash set (\code{std::unordered\_set<uint64\_t>}), costing ${\sim}3.5$\,GB and
running serially (${\sim}200$\,s; an out-of-memory source at $N{=}32$). We
replace it with the identity
\begin{equation}
  s \text{ is a representative} \iff s = \min_{g\in G} g(s).
  \label{eq:repid}
\end{equation}
Site permutations preserve popcount, so every image $g(s)$ is itself a
fixed-$S^z$ state, and the orbit minimum is hit \emph{exactly once} as the loop
sweeps the basis. No deduplication structure is needed:

\begin{lstlisting}[caption={Canonical-representative pass (per-thread, no dedup
set). \code{include/ed/core/streaming\_symmetry.h}},label={lst:pass1}]
for (size_t i = 0; i < num_basis; ++i) {
    const uint64_t basis = basis_states_[i];
    uint64_t rep = basis;
    for (const auto& perm : symmetry_info.max_clique) {
        uint64_t permuted = applyPermutation(basis, perm);
        if (permuted < rep && lookupState(permuted) >= 0) rep = permuted;
    }
    if (basis == rep) thread_reps[tid].push_back(basis);   // canonical
}
\end{lstlisting}

The per-state kernel is a bit-gather, one per group element:
\begin{lstlisting}[caption={\code{applyPermutation} --- output bit $i$ is sourced
from input bit \code{perm[i]}. \code{include/ed/core/basis\_utils.h}}]
inline uint64_t applyPermutation(uint64_t basis, const std::vector<int>& perm) {
    uint64_t result = 0;
    for (size_t i = 0; i < perm.size(); ++i)
        result |= ((basis >> perm[i]) & 1) << i;
    return result;
}
\end{lstlisting}

\paragraph{Complexity.} $\bigO(D\,|G|)$ time, $\bigO(1)$ auxiliary memory. The
loop is OpenMP-parallel (each thread appends to a private vector, then a single
concatenate $+$ \code{std::sort} restores deterministic ascending order).
Measured effect: $200\,\text{s}\rightarrow 6\,\text{s}$ and elimination of the
${\sim}3.5$\,GB set.

\subsection{Sector dimensions without the CSR (``Pass 1.5'')}

Choosing lazy vs.\ eager builds and sizing solver vectors requires each sector's
dimension. Rather than \emph{building} the orbit CSR to discover it, Pass~1.5
walks each orbit, computes only the symmetrized norm of \eqref{eq:symvec}, and
counts survivors, storing a single integer per sector:

\begin{lstlisting}[caption={Norm-only dimension scan; thread-local scratch is
discarded. \code{streaming\_symmetry.h}}]
#pragma omp parallel for schedule(dynamic, 64)
for (size_t oi = 0; oi < num_orbits; ++oi) {
    std::vector<uint64_t> elements; std::vector<Complex> coefficients;
    double norm_sq = 0.0;
    computeOrbitDataFixedSz(unique_orbit_reps_[oi], sector.phase_factors,
                            elements, coefficients, norm_sq);
    if (norm_sq > 1e-10) valid_count.fetch_add(1, std::memory_order_relaxed);
}
symmetrized_block_ham_sizes[sector_idx] = (int)valid_count.load();
\end{lstlisting}

\code{symmetrized\_block\_ham\_sizes[k]} becomes the authoritative dimension,
exposed by \code{getSectorDimension(k)}. This is what lets the GPU sector loop
know \code{dim()} with \emph{zero} orbit materialization. The orbit expansion
goes through the single shared helper \code{compute\_orbit\_for\_state}, so
Pass~1.5, the eager build, and the CSR-free extractor of
\S\ref{sec:repdata} are bit-identical.

\subsection{Lazy LRU-1 materialization}

The eager build materializes every sector's orbit CSR plus reverse index
(${\sim}190$\,GB across $8$ sectors at $N{=}32$). The lazy path materializes
\emph{one} sector at a time at solve/bind time and frees the previous one
(LRU-1), capping instantaneous host RSS at a single sector's footprint
(${\sim}24$\,GiB). This was the prior state of the art but still pays that
$\bigO(D)$ cost while a sector is resident.

\subsection{CSR-free representative data}
\label{sec:repdata}

The decisive step removes the orbit CSR from the GPU path entirely.
\code{getRepSectorData(k)} reproduces the exact survivor set and ordering of the
eager build, but stores only $\bigO(D/|G|)$ data---representatives and
$1/\Norm_\alpha$---using transient scratch for the orbit expansion:

\begin{lstlisting}[caption={CSR-free per-sector source. \code{streaming\_symmetry.h}}]
std::vector<uint64_t> elems; std::vector<Complex> coeffs; double norm_sq = 0.0;
for (uint64_t rep : unique_orbit_reps_) {
    computeOrbitDataFixedSz(rep, sector.phase_factors, elems, coeffs, norm_sq);
    if (norm_sq > 1e-10) {
        out.reps.push_back(rep);
        out.inv_norms.push_back(1.0 / std::sqrt(norm_sq));
    }
}
\end{lstlisting}

It additionally packs the $|G|$ per-sector characters $\chi_k(g)$ and the
flattened group permutations---the only other quantities the device needs. Host
footprint per sector: ${\sim}1.2$\,GiB (\code{reps} $+$ \code{inv\_norms}) versus
${\sim}24$\,GiB (CSR $+$ reverse index).

\paragraph{Lazy host loop.} With this in hand, the production sector loop
(\code{StreamingSymmetryHandle::sector(k)}) hands out a \emph{lazy}
\code{SectorOperator} that (a) knows its dimension up-front (Pass~1.5),
(b) builds the CSR-free record on demand for the GPU bind, and (c) materializes
the orbit CSR \emph{only} if a CPU \code{apply} is ever invoked. A pure-GPU run
therefore never allocates the per-sector orbit CSR on the host.

\section{The symmetrized matrix--vector product}

\subsection{The orbit-CSR SpMV and why it is PCIe-bound}

The orbit-CSR SpMV performs \emph{expand $\rightarrow$ apply $\rightarrow$
project}:
\begin{enumerate}
  \item scatter the ${\sim}75$M reduced vector into the $D$-dim full basis via
        the orbit CSR;
  \item apply $H$ in the full basis;
  \item gather back over the $|G|$ images per orbit via the orbit CSR again.
\end{enumerate}
Steps 1 and 3 each \emph{read} ${\approx}13.44$\,GiB of orbit indices and
coefficients. In lazy mode that CSR lives in host RAM and is streamed over PCIe
(${\sim}50$\,GB/s) every SpMV:
\[
  \frac{13.44\,\text{GiB}\times 2}{50\,\text{GB/s}} \approx 0.54\ \text{s/SpMV},
\]
so the $\div 8$ FLOP win is entirely consumed by host$\leftrightarrow$device
traffic, and the symmetric run is no faster than the unsymmetrized one.

\subsection{On-the-fly representative SpMV: the algebraic collapse}

The new path never expands. It applies $H$ to the \emph{single} representative of
each row and reconstructs the destination index and projection phase
arithmetically. For a connected computational state $s'$ reached from
$r_\alpha$, let $r_\beta=\min_g g(s')$ be its destination representative. The
coefficient of $s'$ in the destination symmetrized vector is
\begin{equation}
  \overline{\beta_{s'}}
   \;=\; \sum_{h:\,h(s')=r_\beta}\overline{\chi_k(h)},
  \qquad
  \text{projection} \;=\; \overline{\beta_{s'}}\,\cdot\,\frac{1}{\Norm_{r_\beta}} .
  \label{eq:proj}
\end{equation}

\paragraph{Derivation.} In \eqref{eq:symvec} the amplitude on $g\,r_\beta$ is
$\chi_k(g)^{*}/\Norm_{r_\beta}$. Writing the set $\{g : g(r_\beta)=s'\}$ as
$\{h^{-1}: h(s')=r_\beta\}$ and using $\chi_k(h^{-1})=\overline{\chi_k(h)}$ for
unit-modulus characters yields \eqref{eq:proj}. The reference's $|G|$-fold orbit
walk and the global $1/|G|$ weight collapse into this single representative
term, so neither appears explicitly. This is implemented in one device
function with no orbit table:

\begin{lstlisting}[caption={\code{DeviceRepSymmetryBasisPolicy::index\_and\_projection}.
\code{include/ed/matvec/device\_basis\_policy.cuh}},label={lst:idxproj}]
__device__ inline std::uint64_t
index_and_projection(std::uint64_t state, cuDoubleComplex& proj_out) const {
    if (__popcll(state) != n_up) return kDeviceNotFound;
    std::uint64_t rb = state;                          // r_b = min_g g(state)
    for (int g = 1; g < group_size; ++g) {
        const std::uint64_t img = apply_perm(state, g);
        if (img < rb) rb = img;
    }
    const int r = ed::gpu::combinadic::rank_state(rb, n_sites, n_up);
    const std::int32_t k = rep_index_of_rank[r];
    if (k < 0) return kDeviceNotFound;
    cuDoubleComplex acc = make_cuDoubleComplex(0.0, 0.0);
    for (int h = 0; h < group_size; ++h)              // sum conj(chi_k(h))
        if (apply_perm(state, h) == rb) {
            const cuDoubleComplex c = characters[h];
            acc.x += cuCreal(c); acc.y -= cuCimag(c); // conjugate
        }
    const double s = inv_norms[k];
    proj_out = make_cuDoubleComplex(acc.x * s, acc.y * s);
    return (std::uint64_t)k;
}
\end{lstlisting}

\code{apply\_perm} is the device twin of \code{applyPermutation}: the entire
group action is regenerated in registers from $|G|\le 8$ permutations and $|G|$
characters---nothing of size $\bigO(D)$ is read or streamed.

\subsection{$\bigO(1)$ reverse lookup: combinadic rank and a dense rank table}

Mapping $r_\beta \mapsto$ orbit index was previously a hash over all $D$ orbit
elements. It is now two $\bigO(1)$ steps.

\paragraph{(a) Combinadic rank.} A fixed-$S^z$ bitmask is converted to its
colexicographic index in $[0,\binom{N}{n_{\uparrow}})$ by reading a Pascal
triangle from \code{\_\_constant\_\_} memory (no global loads):
\begin{equation}
  \mathrm{rank}(s) \;=\; \sum_{b:\,s_b=1}\binom{b}{\,\#\{b'\le b:\,s_{b'}=1\}\,}.
\end{equation}
\begin{lstlisting}[caption={\code{rank\_state}. \code{include/ed/gpu/combinadic.cuh}}]
__device__ __forceinline__
int rank_state(std::uint64_t state, int n_bits, int k) {
    int rank = 0, seen = 0;
    #pragma unroll
    for (int bit = 0; bit < 64; ++bit) {
        if (bit >= n_bits || seen >= k) break;
        if ((state >> bit) & 1ULL) { ++seen; rank += (int)binomial(bit, seen); }
    }
    return rank;
}
\end{lstlisting}

\paragraph{(b) Dense rank table.} \code{rep\_index\_of\_rank} maps
$\mathrm{rank}(\text{rep})\rightarrow$ orbit index, length
$\binom{N}{n_{\uparrow}}$, \code{int32}, with $-1$ for non-representatives. It is
built from \reps{} \emph{only}---no orbit images materialized---and uploaded once
per sector:
\begin{lstlisting}[caption={Reverse table built from reps only.
\code{src/symmetry/streaming\_symmetry\_gpu\_mirror.cu}}]
std::vector<std::int32_t> h_rep_index_of_rank(dim_full_sz, -1);
for (std::size_t i = 0; i < data.reps.size(); ++i) {
    const std::uint64_t r = rank_combination_host(data.reps[i], n_up);
    if (r < dim_full_sz) h_rep_index_of_rank[r] = (std::int32_t)i;
}
\end{lstlisting}
A destination lookup is thus a handful of constant-memory additions
(\code{rank\_state}) followed by a single \code{int32} global load, replacing the
multi-cache-miss hash probe. Device memory for the table is
$\binom{32}{16}\cdot 4\,\text{B}\approx 2.4$\,GiB---negligible on a 40--80\,GB
GPU.

\subsection{The kernel: one thread per representative}

The kernel launches one thread per orbit representative, applies $H$ to
$\reps[i]$ once, and uses $1/\Norm_\alpha=\code{inv\_norms[i]}$ as the source
pre-phase:
\begin{lstlisting}[caption={Rep kernel core. \code{include/ed/matvec/term\_kernels\_gpu.cuh}}]
const double inv_norm_i = basis.inv_norms[i];
const cuDoubleComplex pre_phase = make_cuDoubleComplex(inv_norm_i, 0.0);
process_source_terms<BasisPolicy, Scalar>(
    basis, spin_l, terms, basis.state_of(i), pre_phase, coeff_in, i, out);
\end{lstlisting}
The Hamiltonian term walk \code{process\_source\_terms} is shared verbatim with
the orbit-CSR kernel (extracted into one \code{\_\_device\_\_
\_\_forceinline\_\_} body), so the representative path cannot diverge
numerically from the validated reference. The two policies differ only in
compile-time traits: \code{needs\_orbit\_walk == false} (apply to the single
representative, no CSR sweep) and \code{has\_coeff\_modifier == true} (each
off-diagonal emit routes through \code{index\_and\_projection},
Listing~\ref{lst:idxproj}, for its destination and phase).

\subsection{CPU hot-path lookup (fallback)}

On CPU runs the reverse lookup is a 16\,B/entry sorted array searched with a
branch-free ``monobound'' binary search that prefetches both candidate probe
points one level ahead, compiling the narrow step to a \code{cmov}:
\begin{lstlisting}[caption={\code{SortedUint64Index::find}.
\code{include/ed/core/sorted\_uint64\_index.h}}]
while (len > 1) {
    const std::size_t half = len / 2;
    __builtin_prefetch(a + low + half / 2, 0, 0);
    __builtin_prefetch(a + low + half + half / 2, 0, 0);
    low += (a[low + half] < key) ? half : 0;   // branch-free narrow
    len -= half;
}
\end{lstlisting}
This removes the per-probe mispredict (the hottest op in the CPU symmetry SpMV)
for ${\sim}2\times$; a companion optimization pre-bakes $1/\Norm_\alpha$ so the
inner loop multiplies instead of dividing per emit.

\section{Net accounting}

\begin{table}[h]
\centering
\caption{Per-SpMV cost for one $N{=}32$ fixed-$S^z$ symmetry sector.}
\begin{tabular}{@{}lll@{}}
\toprule
 & Orbit-CSR (lazy) & On-the-fly representative \\
\midrule
Host$\rightarrow$device traffic & $2\times 13.44$\,GiB CSR & two ${\sim}75$M vectors only \\
Device resident                 & ${\sim}30{+}$\,GiB       & ${\sim}6$\,GiB \\
Reverse lookup                  & hash over $D$ elements   & \code{rank\_state} $+$ one \code{int32} load \\
Group action                    & materialized images      & regenerated from $|G|$ perms in registers \\
Host build / sector             & ${\sim}24$\,GiB CSR $+$ index & ${\sim}1.2$\,GiB \reps{}$+$\code{inv\_norms} \\
\bottomrule
\end{tabular}
\end{table}

The orbit-CSR column is PCIe-bound and roughly matches the unsymmetrized cost;
the representative column converts the $|G|$-fold dimension reduction into a real
$|G|$-fold matvec speedup and removes the host-RSS ceiling that previously forced
a $64$\,GB$+$ node. Correctness is pinned bit-for-bit against the CPU
\code{applySymmetrized} reference across $\mathbb{Z}_N$ rings (including a
$|G|{=}8$ complex-character fixed-$S^z$ fixture), and the host-CSR-never-%
materialized invariant on the GPU path is asserted end-to-end through the
production sector loop.

\end{document}
